\documentclass[11pt]{article}

\setlength\textwidth{160mm}
\setlength\textheight{225mm}
\setlength\oddsidemargin{-3mm}
\setlength\headheight{0mm}
\setlength\topmargin{-5mm}

\newtheorem{thm}{Theorem}
\newtheorem{lem}{Lemma}
\newtheorem{cor}{Corollary}
\newtheorem{defn}{Definition}
\newtheorem{fact}{Fact}
\newtheorem{alg}{Algorithm}

\begin{document}

\title{On $\overline{\chi}(G)-\alpha(G)>0$ gap recognition and
$\alpha(G)$-upper bounds}

\author{
Stanislav Busygin\thanks{Email: busygin@ufl.edu. Industrial and Systems 
Engineering Department, University of Florida, 303 Weil Hall, Gainesville, FL 
32611, USA.} \\
Dmitrii V. Pasechnik\thanks{Email: dima@thi.informatik.uni-frankfurt.de. 
Theoretische Informatik, FB15, University of Frankfurt, Robert-Mayer Str. 11-15, 
Postfach 11 19 32, 60054 Frankfurt am Main, Germany. Supported by the DFG Grant 
{\sf SCHN-503/2-1}.}}

\date{}

\maketitle

\begin{abstract}
We show that for a graph $G$ it is $NP$-hard to decide whether
its independence number $\alpha(G)$ equals its clique partition number
$\overline{\chi}(G)$ even when some minimum clique partition of $G$ is
given. This implies that any $\alpha(G)$-upper bound provably better than
$\overline{\chi}(G)$ is $NP$-hard to compute.

To establish this result we use a reduction of the Quasigroup Completion
Problem (QCP, known to be $NP$-complete) to the maximum independent set
problem. A QCP instance is satisfiable if and only if the independence number
$\alpha(G)$ of the graph obtained within the reduction is equal to the number
of holes $h$ in the QCP instance. At the same time, the inequality
$\overline{\chi}(G) \le h$ always holds. Thus, QCP is satisfiable if and
only if $\alpha(G)=\overline{\chi}(G)=h$. Computing the Lov\'asz number
$\vartheta(G)$ we can detect QCP unsatisfiability at least when
$\overline{\chi}(G)<h$. In the other cases QCP reduces to
$\overline{\chi}(G)-\alpha(G)>0$ gap recognition, with one minimum clique 
partition of $G$ known.

{\bf Keywords:} independence number, clique partition number, Lov\'asz number, 
latin square, quasigroup completion problem.
\end{abstract}

\section{Introduction}
Let $G(V,E)$ be a simple undirected graph.
An {\em independent set} of vertices is a subset $S \subseteq V$ such that any
two vertices of $S$ are {\em not} adjacent. The {\em maximum independent set
problem} asks for an independent set of the maximum cardinality. This
cardinality $\alpha(G)$ is called the {\em independence number\/} of the graph, 
and is $NP$-hard to compute \cite{GJ79}. A {\em clique\/} $Q$ is a subset of $V$ 
such that any two vertices of $Q$ are adjacent. The {\em minimum clique 
partition problem} asks for a smallest by cardinality set of cliques
$\{Q_1, \ldots, Q_{\overline{\chi}}\}$ containing every vertex $v \in V$ in exactly one 
of the cliques. The cardinality $\overline{\chi}(G)$ of this set is called the 
{\em clique partition number}. It is equal to the {\em chromatic number} 
$\chi(\overline{G})$ (minimum number of vertex colors
needed to provide different colors for any pair of adjacent vertices) of
the complementary graph. The minimum clique partition problem is also
$NP$-hard \cite{GJ79}. Obviously, the inequality
$\alpha(G) \le \overline{\chi}(G)$ holds as no two vertices of an independent
set can belong to the same clique.

There exists a polynomial-time computable function $\vartheta(G)$
``sandwiched'' between those two $NP$-hard numbers \cite{L79,K94}:
\begin{equation}
\label{eq:sandwich}
\alpha(G) \le \vartheta(G) \le \overline{\chi}(G).
\end{equation}
One simple definition of $\vartheta(G)$ is via minimum of the largest
eigenvalue of so-called {\em feasible matrices} $A=(a_{ij})_{n \times n}$:
\begin{equation}
\label{eq:theta_def}
\vartheta(G) = \min_A \lambda_{\max}(A),
\end{equation}
\[ \mbox{s.t. } a_{ij}=1 \ \mbox{ if } (i,j) \notin E; \quad a_{ij}=a_{ji} \]
(that is, to obtain $\vartheta(G)$ we minimize the largest eigenvalue of
a symmetric matrix having $1$'s on the main diagonal and in all entries
corresponding to non-edges, while the other entries are arbitrary). 
$\vartheta(G)$ is called the {\em Lov\'asz number (or $\vartheta$-function)} of
a graph. It serves as an upper bound for the independence number and as
a lower bound for the clique partition number simultaneously. Besides,
there are increasingly tight sequences of polynomial-time computable
upper bounds for $\alpha(G)$ based on ``lift-and-project'' method \cite{LS91}
and the concept of matrix copositivity \cite{KP01}.

A {\em latin square} is an $n \times n$ matrix filled with integers from $1$
to $n$ so that each number occurs exactly once in any row and in any column. One 
example is $L=(\ell_{ij})_{n \times n}$ such that
\begin{equation}
\label{eq:zn}
\ell_{ij} = ((i+j-2) \ \textrm{mod} \ n) + 1.
\end{equation}
In the {\em Quasigroup Completion Problem} (QCP, a.k.a. latin square completion) 
we are given an $n \times n$ array partially filled with integers from
$\{1 \ldots n\}$ and it is asked whether there is a completion for all $h$ empty 
cells ({\em holes}) such that it gives a latin square. QCP is $NP$-complete 
\cite{C84}. Recently it has been intensively studied, especially from constraint 
programming and boolean satisfiability viewpoints
\cite{GS97,P00,DGKS00,KRAGSS01,GS02,GS02LP,GRS03}.

In this paper we show a reduction of QCP to the maximum independent set
problem. The obtained graph instances obey $\alpha(G) \le h$ and
$\overline{\chi}(G) \le h$ constraints. At that, the original QCP instance
is satisfiable if and only if $\alpha(G)=\overline{\chi}(G)=h$. This allows
us to obtain some results restricting polynomial-time recognition of 
$\overline{\chi}(G)-\alpha(G)>0$ gap and computation of tight upper bounds on 
$\alpha(G)$ under $P \neq NP$ assumption. In fact, unless $P=NP$, it means there 
is no polynomial-time computable upper bound on $\alpha(G)$ provably better than 
$\vartheta(G)$ and, in turn, $\vartheta(G)$ does not bound $\alpha(G)$ provably 
better than $\overline{\chi}(G)$.

\section{Latin square 3D encoding}
The concept of latin square can be also expressed via a three-dimensional
array of 0-1 values. Namely, let a 0-1 variable $x_{ijk}, \ i,j,k \in \{1 \ldots n\}$
denote ``Cell $(i,j)$ is filled with number $k$''. The array of these 
variables determines a latin square if and only if
\begin{equation}
\label{eq:qcp3d}
\left \{ \begin{array}{l}
\forall i,j \ \sum_{k=1}^n x_{ijk} = 1, \\
\forall i,k \ \sum_{j=1}^n x_{ijk} = 1, \\
\forall j,k \ \sum_{i=1}^n x_{ijk} = 1.
\end{array} \right.
\end{equation}
These conditions correspond to maximum independent sets of a graph, whose
vertices are triples $(i,j,k)$ and there is an edge between two of them
if and only if two of their entries coincide. This graph $\Gamma$ is known as 
$H(3,n)$ Hamming graph, see e.g. \cite{BI84}.
\begin{lem}
\label{lem:qcp_3d}
$\alpha(\Gamma)=n^2$. There is a one-to-one correspondence between maximum 
independent sets of $\Gamma$ and $n \times n$ latin squares.
\end{lem}

{\em Proof.\/} First, we prove $\Gamma$ does not have an independent set larger 
in size than $n^2$. Indeed, there are only $n^2$ distinct pairs of two first 
entries $(i,j)$ for the vertices $\{(i,j,k)\}$. Thus, in any vertex subset $X$ 
such that $|X|>n^2$ there is at least one vertex pair $((i,j,k),(p,q,r))$ such 
that $(i,j)=(p,q)$. This vertex pair must be connected by an edge. So, $X$ is 
not an independent set.

Now, consider a latin square $L=(\ell_{ij})_{n \times n}$ and the vertex subset 
$S=\{(i,j,k): \ \ell_{ij}=k\}$. It contains $n^2$ vertices because there are 
$n^2$ distinct $(i,j)$ pairs. Let $(i,j,k) \in S$ and $(p,q,r) \in S$ be two 
distinct vertices. As $i=p$ and $j=q$ would have implied $k=r$ by the definition 
of $S$, this case is not possible. Thus, if $i=p$, then $j \neq q$ and $k \neq r$
as $L$ does not have two equal numbers on the same row. Similarly, if $j=q$, 
then $i \neq p$ and $k \neq r$ as $L$ does not have two equal numbers on the 
same column. Therefore, there are no triples in $S$ with exactly two common 
entries and, hence $L$ defines a maximum independent set of $\Gamma$. As 
(\ref{eq:zn}) provides a latin square for any $n>0$, one obtains 
$\alpha(\Gamma)=n^2$ as claimed.

Conversely, it is easy to see that any $\Gamma$ maximum independent set 
$S=\{(i,j,k)\}, \ |S|=n^2$ defines a latin square $L=(\ell_{ij})_{n \times n}$ 
such that $\ell_{ij}=k$ if and only if $(i,j,k) \in S$. QED.

\vspace{1pc}

To reduce QCP to the maximum independent set problem we will use subgraphs of 
$\Gamma$. Let the QCP input be a matrix $L=(\ell_{ij})_{n \times n}$ such that 
$\ell_{ij}=k \in \{1 \ldots n\}$ if the cell $(i,j)$ is prefilled with $k$, and 
$\ell_{ij}=0$ otherwise. Correspondingly, the number of holes $h$ is the total 
number of entries $(i,j)$ such that $\ell_{ij}=0$. Without loss of generality we 
assume that this input does not immediately violate the latin square 
constraints. That is, $\ell_{ij}=\ell_{iq}>0, \ j \neq q$ or 
$\ell_{ij}=\ell_{pj}>0, \ i \neq p$ cases never occur. Otherwise, the QCP 
instance is trivially unsatisfiable. Define a graph $G(V,E)$ with vertices
\begin{equation}
\label{eq:qcp_v}
V=\{(i,j,k): \ (\ell_{ij}=0) \& (\forall p: \ \ell_{pj} \neq k) \& (\forall q: \ \ell_{iq} \neq k) \}.
\end{equation}
As earlier, put an edge between distinct vertices $(i,j,k)$ and $(p,q,r)$ when 
they have two common entries:
\[ E=\{((i,j,k),(p,q,r)): \ (i=p) \& (j=q) \& (k \neq r) \lor (i=p) \& (j \neq q) \& (k=r) \lor \]
\begin{equation}
\label{eq:qcp_e}
(i \neq p) \& (j=q) \& (k=r)\}.
\end{equation}
In other words, $G(V,E)$ is the subgraph of $\Gamma$ induced by non-neighbors of 
those vertices $(i,j,k)$ for which $\ell_{ij}=k>0$.
\begin{lem}
$\alpha(G) \leq h$. The QCP instance given by the matrix $L$ is satisfiable if 
and only if $\alpha(G)=h$.
\end{lem}

{\em Proof.\/} Let $S_0=\{(i,j,k): \ \ell_{ij}=k>0\}$ be the vertex subset of 
$\Gamma$ corresponding to the partial completion given by $L$. Obviously, 
$|S_0|=n^2-h$. Since the partial completion obeys the latin square constraints, 
$S_0$ is an independent set. Denote by $N^{+}(S_0)$ the closed neighborhood of 
$S_0$, that is, union of $S_0$ with the set of $\Gamma$ vertices adjacent to at 
least one vertex from $S_0$. $G(V,E)$ is obtained by removing $N^{+}(S_0)$ from 
$\Gamma$.

Assume $G(V,E)$ has an independent set $S_1$ of size greater than $h$. Then
$S_0 \cup S_1$ is an independent set of $\Gamma$ having more than $n^2$ vertices, 
contradicting Lemma \ref{lem:qcp_3d}. Hence $\alpha(G) \leq h$.

Let $G(V,E)$ have a maximum independent set $S_1$ of size $h$. Then $S_0 \cup S_1$
is a maximum independent set of $\Gamma$, so $S_1$ determines a correct 
completion of the QCP input to a latin square. Therefore, the given QCP instance 
is satisfiable. Conversely, if the given input matrix $L$ admits a completion to 
a latin square, we can take the maximum independent set $S$ of $\Gamma$ 
corresponding to this latin square and observe that $S \setminus S_0$ is an 
independent set of $G(V,E)$ of size $h$. Therefore, the QCP instance is 
satisfiable if and only if $\alpha(G)=h$. QED.

\vspace{1pc}

Thus, we have described a reduction of QCP to the maximum independent set 
problem on $\Gamma$ subgraphs. In the next section we present the results based 
on clique partition of these subgraphs.

\section{The main results}
\begin{lem}
\label{lem:qcp_chi}
Let $G(V,E)$ be a graph obtained within the QCP reduction to the maximum 
independent set problem. Then $\overline{\chi}(G) \le h$.
\end{lem}

{\em Proof.\/} Let $L=(\ell_{ij})_{n \times n}$ be the QCP input matrix as 
described above. $V$ may include only such vertices $(i,j,k)$ for which 
$\ell_{ij}=0$. We note that all $(i,j,k) \in V$ corresponding to one hole 
$\ell_{ij}=0$ comprise a clique. Hence $V$ is a union of not more than $h$ of 
such cliques. This implies $\overline{\chi}(G) \le h$. QED.

\vspace{1pc}

Therefore, computing the Lov\'asz number $\vartheta(G)$ on the described graphs 
we can efficiently detect QCP unsatisfiability at least when 
$\overline{\chi}(G)<h$. We may say that the inequality $\vartheta(G)<h-\epsilon$ 
for some fixed $0<\epsilon<1$ designates an easily recognizable subclass of 
unsatisfiable QCP instances.
In the other cases, QCP is equivalent to deciding whether 
$\alpha(G)=\overline{\chi}(G)$ provided $\overline{\chi}(G)=h$ and the clique 
partition defined in the proof of Lemma \ref{lem:qcp_chi} is a minimum one. 
Thus, we have deduced the following:
\begin{thm}
\label{thm:main}
For a graph $G$ is it $NP$-hard to decide whether there is a gap between its 
independence and clique partition numbers $\overline{\chi}(G)-\alpha(G)>0$ 
provided some minimum clique partition of $G$ is given.
\end{thm}

We note that currently we are not aware of any graph $G$ obtained within the 
reduction from an unsatisfiable QCP instance for which
$\vartheta(G)=\overline{\chi}(G)=h$.
\begin{cor}
For a graph $G$ is it $NP$-hard to decide whether there is a gap between its 
independence and clique partition numbers $\overline{\chi}(G)-\alpha(G)>0$.
\end{cor}

Though it immediately follows from Theorem \ref{thm:main}, there is also a 
simple direct proof of this fact. Assume we have an oracle answering whether 
$\overline{\chi}(G)-\alpha(G)>0$ for any graph $G$. Define $G_i$ as the graph 
composed of $G$ and $i$ additional mutually independent vertices, each of which 
is connected with every vertex of $G$. Note that $\alpha(G_i)=\max(\alpha(G),i)$ 
and $\overline{\chi}(G_i)=\max(\overline{\chi}(G),i)$. Submit the graphs
$G_i, \ i=0,1,\ldots$ to the oracle until it says $\alpha(G_i)=\overline{\chi}(G_i)$. 
There cannot be more than $\overline{\chi}(G)$ of such queries. Upon 
termination, $\vartheta(G_i)=\overline{\chi}(G)$ since
$\alpha(G_i)=\overline{\chi}(G_i)=\overline{\chi}(G)$, so using the oracle we
can compute $\overline{\chi}(G)$ in polynomial time. (In fact, we have to compute 
$\vartheta(G_i)$ only if the process stops with $i=0$, that is, when 
$\alpha(G)=\overline{\chi}(G)$. Otherwise the terminal value $i$ gives 
$\overline{\chi}(G)$.)

\begin{cor}
Unless $P=NP$, there is no polynomial-time computable upper bound on the 
independence number $\alpha(G)$ provably better than the Lov\'asz number 
$\vartheta(G)$ and, in turn, $\vartheta(G)$ bounds $\alpha(G)$ from above not 
provably better than the clique partition number $\overline{\chi}(G)$.
\end{cor}

Indeed, any such upper bound on $\alpha(G)$ allows for polynomial time 
recognition of $\overline{\chi}(G)-\alpha(G)>0$ gap whenever 
$\overline{\chi}(G)$ is known. According to Theorem \ref{thm:main}, this would 
imply $P=NP$.

\paragraph{Acknowledgements.}
The first author would like to thank G\"unter Stertenbrink for providing him 
with introductory information on QCP and its 3D encoding.

\begin{thebibliography}{13}

\bibitem{GJ79}
  M. Garey and D. Johnson,
  {\em Computers and Intractability: A Guide to the Theory of NP-Completeness\/}
  (Freeman \& Co., 1979).

\bibitem{BI84}
  E. Bannai and T. Ito,
  {\em Algebraic Combinatorics I: Association Schemes\/}
  (Benjamin/Cummings, 1984).
  
\bibitem{L79}
  L. Lov\'asz, On the Shannon capacity of a graph,
  {\em IEEE Trans. Inform. Theory\/} {\bf 25}:1 (1979) 1--7.

\bibitem{C84}
  C. Colbourn,
  The complexity of completing partial latin squares,
  {\em Discrete Applied Mathematics\/}, {\bf 8} (1984) 25--30.

\bibitem{LS91}
  L. Lov\'asz and A. Schrijver,
  Cones of matrices and set-functions and 0-1 optimization,
  {\em SIAM J. Optim.\/} {\bf 1}:2 (1991) 166--190.

\bibitem{K94}
  D.E. Knuth, The Sandwich Theorem,
  {\em Elec. J. Comb.} {\bf 1} (1994).

\bibitem{GS97}
  C. Gomes and B. Selman,
  Problem structure in the presence of perturbations,
  in: {\em Proc. 14th National Conference on Artificial Intelligence
  (AAAI-97)\/} (New Providence, RI, 1997).

\bibitem{P00}
  I. Peterson,
  Completing latin squares,
  {\em Science News\/} {\bf 19}:157 (2000).

\bibitem{DGKS00}
  D. Achlioptas, C. Gomes, H. Kautz, and B. Selman,
  Generating satisfiable problem instances,
  in: {\em Proc. 17th National Conference on Artificial Intelligence
  (AAAI-00)\/} (Austin, TX, 2000).

\bibitem{KRAGSS01}
  H. Kautz, Y. Ruan, D. Achlioptas, C. Gomes, B. Selman, and M. Stickel,
  Balance and filtering in structured satisfiable problems,
  in: {\em Proc. 17th International Conference on Artificial Intelligence
  (IJCAI-2001)\/} (Seattle, WA, 2001).

\bibitem{KP01}
  E. de Klerk and D. Pasechnik,
  Approximation of the stability number of a graph via copositive programming,
  {\em SIAM J. Optim.\/} {\bf 12}:4 (2001) 875--892.

\bibitem{GS02}
  C. Gomes and D. Shmoys,
  Completing quasigroups or latin squares: a structured graph coloring problem,
  in: {\em Proc. Computational Symposium on Graph Coloring and Extensions\/}
  (2002).

\bibitem{GS02LP}
  C. Gomes and D. Shmoys,
  The promise of LP to boost CSP techniques for combinatorial problems,
  in: {\em Proc. 4th International Symposium on Integration of AI and OR
  Techniques in Constraint Programming for Combinatorial Optimization Problems
  (CP-AI-OR'02)\/} (Le Croisic, France, 2002) 291--305.

\bibitem{GRS03}
  C. Gomes, R. Regis, and D. Shmoys,
  An improved approximation for the partial latin square extension problem,
  in: {\em Proc. ACM-SIAM Symposium on Discrete Algorithms (SODA-2003)\/}
  (Baltimore, MD, 2003).

\end{thebibliography}

\end{document}
